\documentclass{report}
\usepackage{amsmath,amssymb}
\setlength{\parindent}{0mm}
\setlength{\parskip}{1em}
\begin{document}
\begin{center}
\rule{6in}{1pt} \
{\large Andreas Stathopoulos \\
{\bf Approximating $Tr(A^{-1})$ with hierarchical probing on regular stencils}}

Department of Computer Science \\ College of William and Mary \\ Williamsburg \\ VA 23187-8795
\\
{\tt andreas@cs.wm.edu}\\
Jesse Laeuchli \\
Konstantinos Orginos\end{center}

The problem of finding the trace of the inverse of a large, sparse
matrix $A$ appears in many applications, including lattice quantum
chromodynamics (LQCD), data mining, and uncertainty quantification.
When the dimension of $A$, $N$, is sufficiently small, the problem is
trivially solved by direct factorization methods.
For large $N$, where factorization of $A$ is not possible, the problem is
surprisingly difficult, and all current approaches resort to a Monte
Carlo method, based on the following:
Given $x \in Z_2^N$, i.e., random vectors of $1,-1$,
$$ E( x^T A^{-1} x) = {\rm Tr}(A^{-1}).$$
A model algorithm generates a list of $s$ vectors, and either solves
a series of linear systems of equations, or computes directly a series
of Gaussian quadratures ($x^TA^{-1}x$).
Because of the rate of Monte Carlo convergence, several hundreds
of systems need to be solved to obtain even a single digit of
accuracy. The need for algorithmic improvements at any level is obvious.

Our driving application is LQCD, where the matrix structure is based on
a four dimensional torus-lattice.
In our previous work, we focused on speeding up the systems with multiple
right hand sides using deflation of a large number of eigenvectors which
are obtained with negligible additional cost to solving the linear systems.
Because of the large number of required right hand sides, we showed
speedups up to an order of magnitude.

In this research we focus on the orthogonal issue of the quality of the
right hand sides. We seek sequences of vectors that are provide better
approximations than random vectors.

Over the last ten years there has been quite some interest in this topic.
Bekas, Kokkiopoulou, and Saad discussed algorithms that use the Hadamard
low disparity vectors that work well for matrices with certain diagonal
structure. They also discussed probing which is our focus below.
Recently, Avron and Toledo did a study of various random vector choices.
Tang and Saad, also discussed a method based on probing.

Probing is a method for obtaining the diagonal (or the trace) of a sparse
matrix. If the graph of the sparse matrix is $k$ colorable,
consider the set of $k$ vectors $X = [x^1, x^2, ..., x^k]$,
where $x^j_i = 1$ if node $i$ has color $j$, and zero otherwise.
It is not hard to see that ${\rm Tr}(A) = {\rm Tr}(X^T AX)$.
For example, in LQCD where the 4-D lattice is 2-colorable, only two
vectors would suffice to find its trace. The problem, however, is that
we need ${\rm Tr}(A^{-1})$, which in general is a dense matrix.

Bekas et al. and Tang et al. exploited the fact that for a large class
of PDE based matrices, the elements of the $A^{-1}$ decay exponentially
away from the non zero elements of $A$. Therefore, the largest elements
of $A^{-1}$ are on the same sparsity pattern as $A$, which for the LQCD
case is $2$-colorable. Thus with 2 vectors we get a rough approximation
of ${\rm Tr}(A^{-1})$. To improve this, we have to consider a ``denser''
approximation of $A^{-1}$, specifically the values of $A^{-1}$ on the
sparsity pattern of $A^2$. This is known as the distance-2 coloring
of the graph of $A$. If this approximation is not satisfactory, we
can continue up to $A^m$, obtain a $k$-coloring of its graph, and use
the $X$ as defined above to obtain the trace. This is the approach
followed by Tang and Saad. However, such approach is wasteful as
all computations performed up to step $m-1$ must be discarded.

In this work, we attempt to avoid this waste; at step $m$ we build
on top of what we have already computed in steps $1,...,m-1$.
However, we are faced with the problem that the vector subspace $X^{(m-1)}$,
computed from the $k^{(m-1)}$ colors needed for $A^{m-1}$, may not be
contained in the subspace $X^{(m)}$ for the next step. Therefore,
previous computations cannot be reused in $X^{(m)T} AX^{(m)}$.
We address this problem by considering a hierarchical graph coloring of $A^m$.
We start with the bipartite, red-black partitioning of $A$ (for LQCD).
Then we consider a subset of the graph of $A^2$ on the red partition only.
We color it with $k_r$ colors.
Then, we color the black partition subset of the $A^2$ with $k_b$ colors.
We can now build $X^{(2)}$ with $k_r+k_b$ vectors, which by construction
also span $X^{(1)}$. Thus, only the complement space needs to be used.
To avoid expensive orthogonalization, we can identify the complement
based on a judicious use of Hadamard vectors.

A concern with the above method is that the number of colors for
any $A^m$ is larger than the minimum (chromatic number of $A^m$).
We are currently investigating how much larger this can be, in
particular for a regular lattice. However, it is unclear whether this is
a real disadvantage, as the sparsity structure of $A^m$ is used only as
an approximation to $A^{-1}$. In fact, more vectors yield more
accurate results.

We will discuss also two variants of this method. First, it is not
necessary to start with a red-black ordering of $A$, as we know that
more vectors will be needed. We can start with a $k$ coloring of $A^m$
(say for $m=2$ or $3$), and then apply the hierarchical graph coloring.
Second, we can stop the hierarchical graph coloring at a certain
distance-$m$, i.e., $A^m$. If more accuracy is needed, each of the
vectors in $X^{(m)}$ can be extended with the regular sequence of
Hadamard vectors, or with random vectors, if we want to remove the bias from
the deterministic approximation.


\end{document}
